\documentclass[12pt]{article}
\usepackage[utf8]{inputenc}
\usepackage[T1]{fontenc}
\usepackage{hyperref}
\title{O'chi'chi' Language Resources for K-12 Curriculum Development}
\author{Compiled by Claude Sonnet 5.5 (claude-sonnet-5-5) for Andr\'as Kornai}
\date{}
\begin{document}
\maketitle

\section{Introduction}
O'chi'chi' (ISO 639-3 xoc; also written \d{O}ch\d{i}ch\d{i}) is reported as a presumably extinct language of Rivers State, Nigeria, in the Etche area (villages reported as Ikwewengwo and Umuebulu). It was first mentioned in print in 2002, and speakers reportedly shifted to the Echie variety of Igbo. Glottolog treats it as a dialect of Obulom and does not recognise it as a separate language, so its status as a distinct language is contested. Search found no teaching materials, word lists beyond the cited encyclopedic entries, or community resources. A heritage and revitalisation curriculum cannot responsibly be built from what is publicly documented.

\section{Primary Educational Resources}
None found.

\section{Linguistic and Reference Materials}
\begin{itemize}
\item Wikipedia, \d{O}ch\d{i}ch\d{i} language: \url{https://en.wikipedia.org/wiki/%E1%BB%8Cch%E1%BB%8Bch%E1%BB%8B_language}
\item Glottolog entry for O'chi'chi' (classified as a dialect of Obulom): \url{https://glottolog.org/resource/languoid/iso/xoc}
\end{itemize}

\section{Religious and Community Resources}
None found.

\section{Digital and Audio Resources}
None found.

\section{Curriculum Development Context}
There is no evidence of living speakers or of documentation sufficient for a curriculum. A reasonable approach would be to contact the Etche community and Nigerian linguists (for example those working on Obulom and Igbo varieties) to establish whether anything is remembered or recorded. Until then, a school unit could address the local linguistic history and the Echie Igbo spoken in the area. No claims about O'chi'chi' vocabulary or grammar should be made without documentation.
\end{document}